% ==========================================
% SIMPLIFIED VERSION - BAB II SECTION 6
% Replace section 6 in Bab II - Studi-Literatur-Revised.tex with this
% ==========================================

\section{Penelitian Basis Utama}
\label{sec:research-foundation}

Berdasarkan tinjauan literatur yang komprehensif, penelitian ini dibangun di atas tiga fondasi metodologis utama yang dipilih karena relevansi langsung dengan karakteristik komunitas Telegram finansial Indonesia yang hierarkis dan memiliki volume pesan tinggi. Bagian ini mengidentifikasi penelitian-penelitian basis yang menjadi referensi utama beserta justifikasi pemilihannya.

\subsection{BERTopic untuk \textit{Topic Modeling} Teks Media Sosial}

Penelitian \textcite{grootendorst2022} tentang BERTopic dipilih sebagai fondasi untuk komponen identifikasi topik diskusi. BERTopic dirancang khusus untuk mengatasi tantangan \textit{topic modeling} pada teks pendek dan informal seperti pesan media sosial. Metodologi ini mengintegrasikan model bahasa berbasis \textit{transformer} untuk \textit{contextual embeddings}, algoritma HDBSCAN untuk \textit{clustering} adaptif, dan prosedur \textit{class-based TF-IDF} untuk ekstraksi representasi topik.

Pemilihan BERTopic dijustifikasi oleh tiga alasan utama. Pertama, superioritas untuk teks pendek yang divalidasi secara empiris oleh \textcite{egger2022} yang membandingkan BERTopic dengan LDA, NMF, dan Top2Vec pada dataset 31.800 \textit{tweet}, mendemonstrasikan keunggulan BERTopic dalam hal koherensi topik dan interpretabilitas. Kedua, kemampuan menghasilkan topik yang interpretatif sangat penting untuk sistem rangkuman yang harus dipahami oleh investor ritel tanpa keahlian teknis NLP. Ketiga, \textit{clustering} berbasis kepadatan dengan HDBSCAN sesuai dengan karakteristik diskusi Telegram yang memiliki topik dengan densitas berbeda-beda.

\subsection{\textit{Weighted Influence System} dari Analisis Sentimen Media Sosial}

Penelitian \textcite{si2014} tentang pemanfaatan sentimen media sosial untuk prediksi pergerakan saham menyediakan fondasi teoretis untuk sistem \textit{weighted influence}. Temuan kunci mereka bahwa "tidak semua opini memiliki bobot yang sama" dan bahwa pendapat dari individu yang lebih berpengaruh atau kredibel dalam jaringan sosial harus diberi bobot lebih tinggi menjadi prinsip desain fundamental sistem rangkuman.

Penelitian \textcite{si2014} menggunakan metrik jaringan sosial seperti \textit{centrality} dan jumlah \textit{followers} untuk menentukan bobot pengaruh. Prinsip bahwa kontribusi dari aktor yang lebih sentral dalam jaringan harus memiliki dampak lebih besar dalam agregasi sentimen diadaptasi untuk konteks komunitas finansial hierarkis. Sistem rangkuman yang diusulkan mengeksploitasi struktur hierarkis eksplisit komunitas Telegram Michael Yeoh untuk memberikan bobot berbeda pada kontribusi dari Michael Yeoh (pemimpin komunitas), administrator, VIP members, dan regular members.

\subsection{Peringkasan Percakapan \textit{Structure-Aware}}

Penelitian \textcite{chen2021} tentang peringkasan percakapan abstraktif yang \textit{structure-aware} menyediakan kerangka metodologis untuk menangani kompleksitas diskusi multi-pihak dalam media sosial. Pendekatan mereka yang memodelkan relasi wacana dan \textit{action triples} ("siapa-melakukan-apa") menggunakan graf terstruktur relevan dengan tantangan meringkas diskusi finansial yang sering kali melibatkan pertukaran argumen, rekomendasi, dan analisis yang kompleks.

Prinsip \textit{structure-awareness} dan konsep \textit{multi-granular decoding} yang dapat menghasilkan rangkuman pada tingkat abstraksi berbeda menginspirasi desain format rangkuman \textit{mixed} yang menggabungkan \textit{bullet points} untuk informasi terstruktur (saham, sinyal, metrik) dengan paragraf naratif untuk konteks dan analisis. Penelitian ini juga menghadapi tantangan tambahan yaitu sintesis informasi dari empat grup dengan karakteristik komunikasi berbeda, yang memerlukan adaptasi pendekatan \textit{structure-aware} untuk konteks multi-grup.

\subsection{Justifikasi Pemilihan dan Integrasi}

Ketiga penelitian ini dipilih karena komplementaritas metodologis mereka dalam menangani aspek-aspek kritis sistem rangkuman untuk komunitas finansial. BERTopic menangani identifikasi topik dari teks informal, \textit{weighted influence system} menangani prioritisasi berdasarkan kredibilitas sumber, dan pendekatan \textit{structure-aware} menangani generasi rangkuman yang koheren dan informatif.

Integrasi ketiga fondasi ini, serta detail implementasi teknis, arsitektur sistem, dan adaptasi spesifik untuk konteks komunitas Telegram finansial Indonesia akan dijelaskan secara komprehensif dalam Bab \ref{chap:desain-konsep-solusi} tentang Desain Konsep Solusi.


\section{Penelitian Basis dan Kerangka Metodologis}
\label{sec:research-foundation}

Berdasarkan tinjauan literatur yang komprehensif, penelitian ini dibangun di atas tiga fondasi metodologis utama yang masing-masing memberikan kontribusi kritis terhadap desain sistem rangkuman otomatis. Bagian ini mengidentifikasi penelitian-penelitian basis yang menjadi referensi utama, menjelaskan aspek-aspek yang diadopsi, serta adaptasi yang dilakukan untuk konteks komunitas finansial Indonesia.

\subsection{Fondasi Pertama: BERTopic untuk \textit{Topic Modeling} Media Sosial}

Penelitian \textcite{grootendorst2022} tentang BERTopic dipilih sebagai fondasi pertama karena metodologinya yang spesifik dirancang untuk mengatasi tantangan \textit{topic modeling} pada teks pendek dan informal seperti media sosial. BERTopic mengintegrasikan empat komponen utama: (1) model bahasa berbasis \textit{transformer} untuk menghasilkan \textit{contextual embeddings}, (2) algoritma UMAP untuk reduksi dimensionalitas, (3) HDBSCAN untuk \textit{clustering} berbasis kepadatan, dan (4) prosedur \textit{class-based TF-IDF} (c-TF-IDF) untuk ekstraksi representasi topik.

\textbf{Aspek yang diadopsi:} Penelitian ini mengadopsi arsitektur \textit{pipeline} BERTopic secara lengkap untuk mengidentifikasi topik diskusi dominan dalam setiap periode rangkuman. Penggunaan \textit{sentence-transformers} untuk menghasilkan \textit{embeddings} semantik dan HDBSCAN untuk \textit{clustering} adaptif sangat sesuai dengan karakteristik pesan Telegram yang bervariasi dalam panjang dan struktur.

\textbf{Adaptasi untuk konteks finansial Indonesia:} Berbeda dengan implementasi standar BERTopic yang menggunakan model \textit{multilingual}, penelitian ini menggunakan model \textit{sentence-transformer} yang di-\textit{fine-tune} untuk bahasa Indonesia (\textit{paraphrase-multilingual-mpnet-base-v2} atau varian Indonesia) untuk meningkatkan kualitas \textit{embeddings}. Selain itu, \textit{stop words} disesuaikan dengan terminologi finansial Indonesia, dan parameter \textit{min\_cluster\_size} dioptimalkan untuk volume pesan komunitas yang tinggi (200+ pesan per jam).

Validasi empiris oleh \textcite{egger2022} yang mendemonstrasikan superioritas BERTopic untuk teks pendek memberikan justifikasi tambahan untuk pemilihan metodologi ini. Kemampuan BERTopic untuk menghasilkan topik yang interpretatif dan koheren sangat penting untuk sistem rangkuman yang harus dipahami oleh investor ritel tanpa keahlian teknis NLP.

\subsection{Fondasi Kedua: \textit{Weighted Influence System} dari Analisis Sentimen Media Sosial}

Penelitian \textcite{si2014} tentang pemanfaatan sentimen media sosial untuk prediksi pergerakan saham menyediakan fondasi teoretis dan metodologis untuk sistem \textit{weighted influence} yang menjadi diferensiator utama penelitian ini. Temuan kunci mereka bahwa "tidak semua opini memiliki bobot yang sama" dan bahwa pendapat dari individu yang lebih berpengaruh atau kredibel dalam jaringan sosial harus diberi bobot lebih tinggi menjadi prinsip desain fundamental.

\textbf{Aspek yang diadopsi:} Konsep \textit{influence-based weighting} diadopsi sebagai mekanisme untuk memprioritaskan informasi berdasarkan kredibilitas sumber. Penelitian \textcite{si2014} menggunakan metrik jaringan sosial seperti \textit{centrality} dan jumlah \textit{followers} untuk menentukan bobot pengaruh. Prinsip bahwa kontribusi dari aktor yang lebih sentral dalam jaringan harus memiliki dampak lebih besar dalam agregasi sentimen menjadi landasan sistem \textit{weighted influence}.

\textbf{Adaptasi untuk struktur hierarkis komunitas:} Berbeda dengan pendekatan \textcite{si2014} yang menggunakan metrik jaringan sosial yang dihitung secara dinamis, penelitian ini mengeksploitasi struktur hierarkis eksplisit komunitas Telegram Michael Yeoh. Sistem \textit{weighting} didefinisikan secara deterministik berdasarkan peran: Michael Yeoh sebagai pemimpin dan ahli utama (bobot 4×), administrator yang memiliki kredibilitas tinggi (bobot 3×), VIP members yang merupakan investor berpengalaman (bobot 2×), dan regular members (bobot 1×). Pendekatan ini lebih sederhana secara komputasional namun tetap menangkap hierarki kredibilitas yang relevan dengan konteks komunitas finansial.

Implementasi \textit{weighted influence} mempengaruhi tiga aspek sistem: (1) ekstraksi saham yang paling relevan dengan memberikan bobot lebih tinggi pada ticker yang disebutkan oleh sumber kredibel, (2) agregasi sentimen dengan menghitung rata-rata tertimbang sentimen berdasarkan bobot pengirim, dan (3) prioritisasi pesan untuk peringkasan dengan memberikan skor relevansi yang lebih tinggi pada kontribusi dari sumber berpengaruh.

\subsection{Fondasi Ketiga: Peringkasan Percakapan \textit{Structure-Aware}}

Penelitian \textcite{chen2021} tentang peringkasan percakapan abstraktif yang \textit{structure-aware} menyediakan kerangka metodologis untuk menangani kompleksitas diskusi multi-pihak dalam media sosial. Pendekatan mereka yang memodelkan relasi wacana dan \textit{action triples} ("siapa-melakukan-apa") menggunakan graf terstruktur sangat relevan dengan tantangan meringkas diskusi finansial yang sering kali melibatkan pertukaran argumen, rekomendasi, dan analisis yang kompleks.

\textbf{Aspek yang diadopsi:} Prinsip \textit{structure-awareness} diadopsi dalam desain \textit{prompt engineering} untuk LLM. Sistem dirancang untuk tidak hanya mengagregasi konten semantik, tetapi juga mempertahankan struktur argumentatif dan relasi antar pesan. Konsep \textit{multi-granular decoding} yang dapat menghasilkan rangkuman pada tingkat abstraksi berbeda menginspirasi desain format rangkuman \textit{mixed} yang menggabungkan \textit{bullet points} untuk informasi terstruktur (saham, sinyal, metrik) dengan paragraf naratif untuk konteks dan analisis.

\textbf{Adaptasi untuk sintesis lintas-grup:} Penelitian \textcite{chen2021} fokus pada peringkasan satu percakapan, sementara penelitian ini menghadapi tantangan tambahan: sintesis informasi dari empat grup dengan karakteristik komunikasi berbeda (broadcast di Cuap Cuap, sinyal terstruktur di Swing Plan, diskusi formal di Advanced, debat kasual di PIRANHA). Adaptasi dilakukan dengan menerapkan struktur graf pada level yang lebih tinggi: setiap grup dipandang sebagai satu node dalam graf meta-komunikasi, dan relasi antar-grup dimodelkan sebagai edge yang menghubungkan informasi terkait (misalnya, sinyal dari Swing Plan yang diimplementasikan dan didiskusikan di Advanced).

Mekanisme \textit{cross-group synthesis} diimplementasikan melalui \textit{prompt engineering} yang eksplisit: LLM diberi instruksi untuk mengidentifikasi hubungan semantik antar informasi dari grup berbeda dan mensintesisnya menjadi narasi koheren yang memberikan pandangan holistik. Misalnya, analisis makroekonomi dari Cuap Cuap dapat dihubungkan dengan sinyal konkret dari Swing Plan dan diskusi implementasi dari Advanced dalam satu rangkuman terintegrasi.

\subsection{Integrasi Tiga Fondasi dalam Arsitektur Sistem}

Ketiga fondasi metodologis ini tidak beroperasi secara terpisah, melainkan diintegrasikan dalam arsitektur \textit{pipeline} yang kohesif. Alur kerja sistem dapat dipahami sebagai berikut:

\begin{enumerate}
\item \textbf{Fase Ekstraksi dan Analisis (BERTopic + \textit{Weighted Influence}):} Pesan dari empat grup dikumpulkan dalam jendela waktu satu jam. BERTopic diterapkan untuk mengidentifikasi topik dominan. Secara paralel, NER mengekstrak ticker saham dan analisis emosi mengklasifikasikan sentimen psikologis. Sistem \textit{weighted influence} memberikan bobot pada setiap pesan berdasarkan peran pengirim.

\item \textbf{Fase Agregasi dan Prioritisasi (\textit{Weighted Influence}):} Ticker saham yang terdeteksi diagregasi dengan bobot berdasarkan kredibilitas sumber yang menyebutkannya. Sentimen dan emosi diagregasi menggunakan rata-rata tertimbang. Pesan diprioritaskan untuk konteks LLM berdasarkan kombinasi relevansi topik, bobot pengirim, dan kebaruan informasi.

\item \textbf{Fase Sintesis dan Generasi (\textit{Structure-Aware Summarization}):} LLM menerima input terstruktur yang mencakup: (1) topik dominan dari BERTopic, (2) ticker teragregasi dengan sentimen tertimbang, (3) subset pesan prioritas dari keempat grup, dan (4) instruksi eksplisit untuk sintesis lintas-grup. Output yang dihasilkan adalah rangkuman \textit{mixed-format} yang mempertahankan struktur informasi sambil memberikan narasi koheren.
\end{enumerate}

Integrasi ini menghasilkan sistem yang lebih dari sekadar jumlah komponennya. \textit{Weighted influence system} meningkatkan kualitas input untuk BERTopic dengan memberikan bobot lebih tinggi pada pesan dari sumber kredibel, yang pada gilirannya menghasilkan topik yang lebih representatif dari diskusi substantif daripada noise. Topik yang teridentifikasi kemudian menginformasikan prioritisasi pesan untuk konteks LLM, memastikan bahwa generasi rangkuman fokus pada informasi yang paling relevan. Akhirnya, pendekatan \textit{structure-aware} memastikan bahwa rangkuman yang dihasilkan bukan hanya agregasi informasi, tetapi sintesis yang koheren dengan struktur argumentatif yang jelas.

\subsection{Kontribusi Teoretis dan Praktis}

Dengan membangun pada tiga fondasi ini dan mengadaptasinya untuk konteks komunitas finansial Indonesia yang hierarkis dan multi-grup, penelitian ini memberikan beberapa kontribusi:

\textbf{Kontribusi teoretis:} Penelitian ini memperluas teori \textit{information overload} dan \textit{Cognitive Load Theory} ke konteks baru: komunitas finansial berjenjang berbasis Telegram. Integrasi \textit{weighted influence} dengan \textit{topic modeling} dan peringkasan percakapan dalam satu \textit{pipeline} kohesif merupakan pendekatan novel yang belum banyak dieksplorasi dalam literatur.

\textbf{Kontribusi metodologis:} Adaptasi BERTopic untuk bahasa Indonesia dengan optimasi parameter untuk teks finansial, implementasi \textit{weighted influence} berbasis struktur hierarkis eksplisit, dan desain \textit{prompt engineering} untuk sintesis lintas-grup menyediakan \textit{blueprint} metodologis yang dapat diadaptasi untuk komunitas online serupa di domain lain.

\textbf{Kontribusi praktis:} Sistem yang dikembangkan menangani masalah \textit{real-world} dengan dampak langsung bagi 22.000+ anggota komunitas Michael Yeoh. Validasi melalui deployment aktif dan \textit{feedback} pengguna menyediakan bukti empiris tentang efektivitas pendekatan dalam mereduksi \textit{information overload} dan meningkatkan efisiensi akses informasi.

\textbf{Kontribusi untuk NLP Indonesia:} Penelitian ini menyediakan implementasi \textit{end-to-end} sistem NLP berbasis AI untuk domain finansial Indonesia yang masih \textit{underrepresented} dalam literatur. Integrasi model NER bahasa Indonesia, klasifikasi emosi, dan generasi teks dengan LLM dalam konteks aplikasi \textit{real-world} memberikan \textit{case study} berharga untuk pengembangan aplikasi NLP Indonesia di domain spesifik lainnya.
